\section{Metodologia}
\label{chap:metodologia}

Para o desenvolvimento deste trabalho, implementamos e validamos quatro variações algorítmicas para o cálculo da matriz de covariância ($\mathbf{C}$), fundamentais para a caracterização estatística das classes no dataset \textit{Wall-Following Robot Navigation}. Seja $\mathbf{X} \in \mathbb{R}^{p \times N}$ a matriz de dados, onde $p$ é o número de sensores e $N$ o número de amostras. Seja $\mathbf{m} = \frac{1}{N} \sum_{j=1}^N \mathbf{x}_j$ o vetor de médias aritmética de cada sensor. A matriz de covariância sumariza as relações lineares entre os atributos, sendo definida pela esperança dos produtos externos dos desvios em relação à média.

O primeiro método (\texttt{mcovar1}) utiliza um laço explícito para acumular os produtos externos dos desvios de cada amostra: $\mathbf{C}_1 = \frac{1}{N} \sum_{j=1}^N (\mathbf{x}_j - \mathbf{m})(\mathbf{x}_j - \mathbf{m})^T$. Embora intuitiva, esta abordagem apresenta alto \textit{overhead} no interpretador Python devido à iteração sobre as $N$ amostras. O segundo método (\texttt{mcovar2}) otimiza este cálculo via vetorização em bloco, utilizando a matriz de dados centralizada $\tilde{\mathbf{X}} = \mathbf{X} - \mathbf{m} \mathbf{1}^T$ para computar $\mathbf{C}_2 = \frac{1}{N} \tilde{\mathbf{X}} \tilde{\mathbf{X}}^T$. Esta abordagem aproveita rotinas BLAS (\textit{Basic Linear Algebra Subprograms}) de alto desempenho, que utilizam paralelismo em nível de instrução (SIMD) e otimizações de cache da CPU.

As outras duas abordagens baseiam-se na relação entre correlação e covariância. O terceiro método (\texttt{mcovar3}) calcula a matriz de correlação $\mathbf{R}$ via laço e subtrai o produto externo das médias ($\mathbf{C}_3 = \mathbf{R} - \mathbf{m} \mathbf{m}^T$), enquanto o quarto método (\texttt{mcovar4}) realiza esta operação de forma puramente matricial: $\mathbf{C}_4 = \frac{1}{N} \mathbf{X} \mathbf{X}^T - \mathbf{m} \mathbf{m}^T$. O protocolo experimental submeteu estas implementações a um \textit{benchmark} de 100 rodadas independentes, comparando-as com a função nativa do NumPy em termos de latência média e precisão numérica, quantificada pela norma de Frobenius do erro residual em relação à referência.

A estabilidade numérica foi investigada por meio da análise de invertibilidade. Em problemas de robótica com sensores redundantes, as matrizes de covariância frequentemente apresentam mal-condicionamento crítico, onde o número de condicionamento $\kappa(\mathbf{C})$ é extremamente elevado. Nestes casos, a matriz torna-se singular para fins computacionais, impossibilitando sua inversão para o cálculo de densidades Gaussianas. Para mitigar este problema, empregamos a regularização de Tikhonov ($\mathbf{C}_{reg} = \mathbf{C} + \lambda \mathbf{I}$), com $\lambda = 0{,}1$, garantindo que a matriz resultante possua posto completo e seja estável para operações subsequentes.
